\chapter{A Comunidade Humana}

A vocação da humanidade é manifestar a imagem de Deus e ser transformada à imagem do Filho único do Pai. Esta vocação reveste-se de uma forma pessoal, pois cada um é chamado a entrar na bem-aventurança divina, mas diz também respeito ao conjunto da comunidade humana. Nos parágrafos 1877 a 1948, o Catecismo da Igreja Católica examina como a vida em sociedade não é um acidente da história, mas uma exigência inscrita na própria natureza do homem criado à imagem de um Deus que é comunhão. Este capítulo percorre três grandes eixos: a relação entre a pessoa e a sociedade, as condições para uma participação justa na vida social e as exigências da justiça social. Em cada um deles ressoa a mesma convicção: a vocação do homem a \textbf{fazer o bem} não se realiza em isolamento, mas no seio de uma comunidade que deve ser construída sobre a verdade, a justiça e a caridade.

\section{A pessoa e a sociedade (1877--1896)}

Todos os homens são chamados ao mesmo fim, que é o próprio Deus. Existe uma certa semelhança entre a unidade das pessoas divinas e a fraternidade que os homens devem instaurar entre si, na verdade e no amor; o amor ao próximo é, por isso, inseparável do amor a Deus. A pessoa humana tem necessidade da vida social: esta não é algo de acessório à sua natureza, mas uma exigência dela; graças ao contacto com os outros, ao serviço mútuo e ao diálogo com os seus irmãos, o homem desenvolve as suas capacidades e assim responde à sua vocação. A sociedade é um conjunto de pessoas ligadas de modo orgânico por um princípio de unidade que ultrapassa cada uma delas; através dela, cada homem é constituído ``herdeiro'', recebendo ``talentos'' que enriquecem a sua identidade e cujos frutos deve desenvolver. Com toda a razão, cada um é devedor de dedicação às comunidades de que faz parte e de respeito às autoridades encarregadas do \textbf{bem} comum.

Cada comunidade define-se pelo fim a que tende, mas ``a pessoa humana é e deve ser o princípio, o sujeito e o fim de todas as instituições sociais'' (GS 25). Certas sociedades, como a família e a comunidade civil, correspondem mais imediatamente à natureza do homem e são-lhe necessárias. Deve fomentar-se a criação de associações e instituições de livre iniciativa --- a ``socialização'' que desenvolve as qualidades da pessoa, particularmente o sentido de iniciativa e de responsabilidade. Mas a socialização também oferece perigos: uma intervenção exagerada do Estado pode ameaçar a liberdade e as iniciativas pessoais. A doutrina da Igreja elaborou o princípio da subsidiariedade: ``uma sociedade de ordem superior não deve interferir na vida interna duma sociedade de ordem inferior, privando-a das suas competências, mas deve apoiá-la e ajudá-la a coordenar a sua acção com a dos demais componentes sociais, com vista ao \textbf{bem} comum'' (João Paulo II, \textit{Centesimus annus}, 48). Este princípio opõe-se a todas as formas de coletivismo e visa instaurar uma verdadeira ordem que respeite a liberdade humana a todos os níveis.

A sociedade é indispensável à realização da vocação humana, mas deve para isso ser respeitada a justa hierarquia dos valores, que subordina as dimensões físicas e instintivas às dimensões interiores e espirituais. A convivência humana é, antes de mais, um facto de ordem espiritual: comunicação de conhecimentos à luz da verdade, exercício de direitos e cumprimento de deveres, partilha dos \textbf{bens} do espírito e aspiração a uma mútua assimilação de valores. A inversão dos meios e dos fins --- dar valor de fim último ao que não é mais que meio, ou considerar as pessoas como puros meios --- gera estruturas injustas que tornam ``árduo e praticamente impossível'' um procedimento cristão conforme aos mandamentos divinos. É preciso, portanto, apelar para as capacidades espirituais e morais da pessoa e para a exigência permanente da sua conversão interior para se conseguirem mudanças sociais que estejam realmente ao seu serviço, impondo a obrigação de introduzir nas instituições as correções convenientes, para que elas se conformem com as normas da justiça e favoreçam o \textbf{bem}, em vez de se lhe oporem. Sem a ajuda da graça, os homens não seriam capazes de descobrir o caminho estreito entre a cobardia que cede ao mal e a violência que o agrava; é o caminho da caridade, que constitui o maior mandamento social e inspira uma vida de entrega ao próximo.

\section{A participação na vida social (1897--1927)}

``A sociedade humana não estará bem constituída nem será fecunda, se a ela não presidir uma autoridade legítima que salvaguarde as instituições e dedique o necessário trabalho e esforço ao \textbf{bem} comum'' (João XXIII, \textit{Pacem in terris}, 46). Toda a comunidade humana tem necessidade de uma autoridade que a governe, fundada na natureza humana e necessária para a unidade da comunidade civil; o seu papel consiste em assegurar, quanto possível, o \textbf{bem} comum da sociedade. A autoridade exigida pela ordem moral emana de Deus: ``Não há autoridade que não tenha sido constituída por Deus'' (Rm 13,1). O dever de obediência impõe a todos a obrigação de tributar à autoridade as honras que lhe são devidas; São Clemente de Roma redigiu para isso a mais antiga oração da Igreja pela autoridade política: ``Dirigi, Senhor, o seu conselho segundo o que é \textbf{bem}, segundo o que é agradável aos vossos olhos, para que, exercendo com piedade, na paz e na mansidão, o poder que lhes destes, Vos encontrem propício.''

A determinação dos regimes políticos e a designação dos seus dirigentes devem ser deixadas à livre vontade dos cidadãos. A diversidade dos regimes políticos é moralmente admissível, desde que concorram para o \textbf{bem} legítimo da comunidade que os adota; os regimes cuja natureza for contrária à lei natural, à ordem pública e aos direitos fundamentais das pessoas não podem promover o \textbf{bem} comum. A autoridade não recebe de si mesma a legitimidade moral: ``A legislação humana só se reveste do carácter de lei, na medida em que se conforma com a justa razão; daí ser evidente que ela recebe todo o seu vigor da lei eterna'' (São Tomás de Aquino). A autoridade só é exercida legitimamente quando procura o \textbf{bem} comum e emprega para o atingir meios moralmente lícitos. ``É preferível que todo o poder seja equilibrado por outros poderes e outras competências que o mantenham no seu justo limite. Este é o princípio do `Estado de direito', no qual é soberana a lei, e não a vontade arbitrária dos homens'' (João Paulo II, \textit{Centesimus annus}, 44).

Em conformidade com a natureza social do homem, o \textbf{bem} de cada um está necessariamente relacionado com o \textbf{bem} comum, que não pode definir-se senão em referência à pessoa humana. Por \textbf{bem} comum deve entender-se ``o conjunto das condições sociais que permitem, tanto aos grupos como a cada um dos seus membros, atingir a sua perfeição, do modo mais completo e adequado'' (GS 26). O \textbf{bem} comum interessa à vida de todos e inclui três elementos essenciais: o respeito da pessoa como tal e dos seus direitos fundamentais inalienáveis; o \textbf{bem}-estar social e o desenvolvimento da sociedade, para que cada um possa aceder ao que lhe é necessário para uma vida verdadeiramente humana; e a paz, isto é, a permanência e a segurança de uma ordem justa. O \textbf{bem} comum está sempre orientado para o progresso das pessoas: ``A ordem das coisas deve estar subordinada à ordem das pessoas, e não o inverso'' (GS 26). A participação de todos na promoção do \textbf{bem} comum implica uma conversão incessantemente renovada dos parceiros sociais, pois incumbe àqueles que exercem cargos de autoridade garantir os valores que atraem a confiança dos membros do grupo e os incitam a colocar-se ao serviço dos seus semelhantes.

\section{A justiça social (1928--1948)}

A sociedade garante a justiça social quando realiza as condições que permitem às associações e aos indivíduos obterem o que lhes é devido, segundo a sua natureza e vocação. A justiça social está ligada ao \textbf{bem} comum e ao exercício da autoridade. A justiça social só pode alcançar-se no respeito da dignidade transcendente do homem: a defesa e a promoção da dignidade da pessoa humana ``foram-nos confiadas pelo Criador, tarefa a que estão rigorosa e responsavelmente obrigados os homens e as mulheres em todas as conjunturas da história'' (João Paulo II, \textit{Sollicitudo rei socialis}, 47). O respeito pela pessoa humana implica o dos direitos que dimanam da sua dignidade de criatura; esses direitos são anteriores à sociedade e impõem-se-lhe, estando na base da legitimidade moral de qualquer autoridade. O respeito pela pessoa humana passa ainda pelo respeito do princípio: ``Que cada um considere o seu próximo, sem qualquer excepção, como `outro ele mesmo', e zele, antes de mais, pela sua existência e pelos meios que lhe são necessários para viver dignamente'' (GS 27). A caridade vê em cada homem um ``próximo'', um irmão, e o dever de nos fazermos o próximo do outro é tanto mais premente quanto esse outro for mais indefeso: ``Quantas vezes o fizestes a um dos meus irmãos mais pequeninos, a Mim o fizestes'' (Mt 25,40).

Criados à imagem do Deus único, dotados de idêntica alma racional, todos os homens têm a mesma natureza e a mesma origem; resgatados pelo sacrifício de Cristo, todos são chamados a participar da mesma bem-aventurança divina e gozam, portanto, de igual dignidade. A igualdade entre os homens assenta essencialmente na sua dignidade pessoal e nos direitos que dela dimanam; toda a forma de discriminação relativamente aos direitos fundamentais da pessoa, por razão do sexo, raça, cor, condição social, língua ou religião, ``deve ser ultrapassada e eliminada como contrária ao desígnio de Deus'' (GS 29). Existem, porém, diferenças entre os homens: fazem parte do plano de Deus, que quer que cada um receba de outrem aquilo de que precisa, e que os que dispõem de talentos particulares comuniquem os seus benefícios a quem deles precisa. As diferenças devem estimular e muitas vezes obrigam as pessoas à magnanimidade, à benevolência e à partilha. Mas também existem desigualdades iníquas que ferem milhões de homens e de mulheres; essas estão em contradição frontal com o Evangelho: ``As excessivas desigualdades económicas e sociais entre os membros ou povos da única família humana provocam escândalo e são obstáculo à justiça social, à equidade, à dignidade da pessoa humana e, finalmente, à paz social e internacional'' (GS 29).

O princípio da solidariedade, também enunciado sob o nome de ``amizade'' ou de ``caridade social'', é uma exigência direta da fraternidade humana e cristã. A solidariedade manifesta-se, em primeiro lugar, na repartição dos \textbf{bens} e na remuneração do trabalho; implica também o esforço por uma ordem social mais justa, em que as tensões possam ser resolvidas melhor e os conflitos encontrem uma saída negociada. Os problemas socioeconômicos só podem ser resolvidos com a ajuda de todas as formas de solidariedade: entre os pobres, dos ricos com os pobres, entre os trabalhadores, entre empresários e empregados, entre as nações e entre os povos. A virtude da solidariedade vai além dos \textbf{bens} materiais: ao difundir os \textbf{bens} espirituais da fé, a Igreja favoreceu também o desenvolvimento dos \textbf{bens} temporais, abrindo-lhes muitas vezes novos caminhos; assim se verificou, ao longo dos séculos, a Palavra do Senhor: ``Procurai primeiro o Reino de Deus e a sua justiça, e tudo o mais vos será dado por acréscimo'' (Mt 6,33). A solidariedade cristã lembra que não existe solução para a questão social fora do Evangelho, e que a caridade --- que respeita o outro e os seus direitos, exige a prática da justiça e inspira uma vida de entrega --- continua a ser o maior mandamento social, a força que desde há dois mil anos leva as almas ``até ao heroísmo caridoso'' em favor de condições sociais capazes de tornar possível a todos uma vida digna do homem e do cristão.
